\documentclass[a4paper]{article}
\usepackage[pdftex]{graphicx}
\usepackage[utf8]{inputenc}
\usepackage{enumerate}
\usepackage{siunitx}
\sisetup{locale=DE} 
\usepackage{href-ul}
\hypersetup{
	colorlinks=true,
	linkcolor=blue,
	urlcolor=blue}
\usepackage{icomma}
\usepackage{amssymb}
\usepackage{eurosym}
\usepackage{geometry}
\geometry{a4paper, top=15mm, left=15mm, right=15mm, bottom=15mm,
	headsep=10mm, footskip=12mm}
%Source: img_0107

% Start the document
\begin{document}
%Titel
{\bf Schulaufgabe aus der Physik }\\
\begin{enumerate}[1.]

%Aufgabe 1
{\bf \item Aufgabe}
\begin{enumerate}[a)]
\item Beschreibe einen Versuch, durch den wir im Unterricht zeigen können, dass sich feste Körper bei Erwärmung ausdehnen.
\item Beschreibe mithilfe des Teilchenmodells, wie sich dabei die innere Energie dieses Körpers verändert
\item Um wieviel ändert sich die Länge einer 30 m langen Eisenbahnschiene, wenn sie einer Temperaturschwankung zwischen $\SI{-25}{\celsius}$ und $\SI{+40}{\celsius}$ ausgesetzt ist. --- Ein 1 m langer Stahlstab ändert bei Erwärmung um 1 K seine Länge um 0,012 mm.
\end{enumerate}

\begin{minipage}{0.65\textwidth}
%Aufgabe 2
{\bf \item Ein Trog, der mit 7,0 Liter Wasser gefüllt ist, steht auf einem 50 m hohen Turm (siehe Skizze).} In dem Trog befindet sich ein Rührwerk, das durch ein langsam sinkendes Massestück G von 100 kg in Gang gehalten wird. Das Wasser im Trog gerät in heftige Bewegung und beruhigt sich dann wieder vollständig. 

\begin{enumerate}[a)]
\item Welche Energieumwandlungen treten bei diesem Versuch auf? --- Kurze Darstellung in der richtigen Reihenfolge!
\item Welche Temperaturdifferenz erfährt die Wassermenge aufgrund dieses Vorgangs, wenn etwa 60\% der zur Verfügung stehenden Energie an das Wasser abgegeben wird?[$c_w=\SI{4,18}{\kilo\joule\over \kilogram\kelvin}$]
\end{enumerate}
\end{minipage}
\hfill
\begin{minipage}{0.3\textwidth}
\includegraphics[width=4 cm]{turuehr107.pdf}
\end{minipage}

%Aufgabe 3
{\bf \item Aufgabe}
\begin{enumerate}[a)]
\item Herr Heiß bereitet sich ein Bad mit 150 Liter Wasser der Temperatur 
$\vartheta_2= \SI{38}{\celsius}$. Das zu erwärmende Wasser hatte $\vartheta_1= \SI{14}{\celsius}$.
Wie viel kostet dieses Wannenbad, wenn für $\SI{1000}{\kilo\joule}$ etwa 4 Cent zu bezahlen sind?
\item Wie viel Energie würde Herr Heiß beim Duschen sparen, wenn die Temperaturwerte des Wassers gleich bleiben und zum Duschen etwa 25 Liter Wasser verbraucht werden?
\end{enumerate}

{\bf \item Die Familie Berger will sich zu später Stunde einen Kaffee kochen.} Die Ma\-schine fasst 1,2 Liter Wasser. 
\begin{enumerate}[a)]
\item Wie viel Energie brauchen sie für $\SI{80}{\celsius}$ heißes Kaffeewasser, wenn das Wasser anfangs  $\SI{15}{\celsius}$ besitzt?
\item Wie lange braucht die Maschine dafür, wenn sie die Leistung von 800 W liefern kann?
\end{enumerate} 

\end{enumerate}

\fbox{
	\begin{minipage}{0.5\textwidth}
		Zur Lösung bitte \href{https://www.okuyakl.de/phys/p9thbL107/ll107.pdf}{hier klicken} oder den QR-Code scannen.\\
	Weitere Arbeitsblätter gibt es unter 
	
	\href{https://www.okuyakl.de}{www.okuyakl.de}
	\end{minipage}
	\hfill
	\begin{minipage}{0.4\textwidth}
		\includegraphics[width=1.5 cm]{../../viecher/zwe03}
		\includegraphics[width=3 cm]{qr107}
		\includegraphics[width=2 cm]{../../viecher/afanticon1}
		
	\end{minipage}}
\end{document}
